\chapter{Anatomy of an Ada Application}
\label{ch:anatomy}

An Ada program for the ESP32-S3 is not a flat blob. It is a stack of layers, each
written in Ada (with a thin sliver of Xtensa assembly at the very bottom). Knowing
the layers explains what links into your binary, what runs before \texttt{Main},
and where each feature comes from.

\section{The layers, top to bottom}

\begin{description}[style=nextline]
\item[Your application.] One or more Ada packages and a main procedure
  (\texttt{Main}). It \texttt{with}s the peripheral HAL (\texttt{ESP32S3.*}) and,
  if it uses tasks or protected objects, those language features directly.
\item[The HAL (\texttt{libs/esp32s3\_hal}).] Hand-written driver packages over a
  generated register layer (\texttt{ESP32S3\_Registers.*}). Optional --- a program
  can poke registers itself --- but it is where all the device knowledge lives.
\item[The GNAT binder output.] \texttt{gnatbind} generates a package
  (\texttt{b\_\_main}) that elaborates every unit in the right order and calls
  \texttt{Main}. Its \texttt{adainit}/\texttt{adafinal} bracket your program.
\item[The Ada runtime (RTS).] The \texttt{System.*} and \texttt{Ada.*} bodies:
  secondary stack, exceptions (if enabled), \texttt{Ada.Text\_IO} over the
  USB-serial console, and --- crucially --- the tasking kernel (GNARL/GNARLI):
  tasks, protected objects, delays, scheduling.
\item[Board support.] The silicon-specific bottom of the runtime: the Xtensa
  windowed \emph{context switch}, \emph{interrupt entry/exit} vectors, the
  \emph{clock/alarm tick}, SMP scheduling and the inter-core interrupt. This is
  Ada plus a little assembly (\texttt{start.S}, \texttt{highint5.S},
  \texttt{context.S}). This is the layer an off-the-shelf RTOS would otherwise
  provide --- here it is ours.
\item[The 2nd-stage bootloader.] Before any of the above runs, our own minimal
  loader (a ZFP-Ada \texttt{boot\_main} + an assembly prologue + a thin C PSRAM
  shim) maps the flash cache/MMU and brings up the octal PSRAM, then jumps to the
  application image. It replaces the chip vendor's stock second-stage bootloader.
\item[The mask ROM.] The chip's on-chip ROM loads the 2nd-stage bootloader from
  flash at reset.
\end{description}

\begin{figure}[ht]
\centering
\begin{tikzpicture}[node distance=2mm]
\node[layer, fill=blue!9] (app)  {Your application --- \texttt{Main} + packages};
\node[layer, below=of app]  (hal)  {Peripheral HAL: \texttt{ESP32S3.*} over generated \texttt{ESP32S3\_Registers.*}};
\node[layer, below=of hal]  (bind) {GNAT binder --- \texttt{adainit} / \texttt{adafinal}, calls \texttt{Main}};
\node[layer, below=of bind] (rts)  {Ada runtime: \texttt{System.*}, \texttt{Ada.*}, tasking kernel (GNARL)};
\node[layer, below=of rts]  (bsp)  {Board support: context switch, interrupt vectors, tick, SMP\, (\texttt{*.S})};
\node[layer, below=of bsp]  (boot) {2nd-stage bootloader --- cache/MMU, octal PSRAM};
\node[layer, below=of boot, fill=gray!12] (rom) {Mask ROM --- loads the bootloader at reset};
\end{tikzpicture}
\caption{The layers of an Ada application, from your code down to the chip's ROM.
Everything above the ROM is Ada (with a little Xtensa assembly in board support).}
\label{fig:layers}
\end{figure}

\section{What runs before \texttt{Main}}

At power-on: ROM $\rightarrow$ our 2nd-stage bootloader (cache/MMU + PSRAM)
$\rightarrow$ \texttt{start.S} (selects the 240\,MHz PLL, sets up the stack and
the interrupt vectors, and on the SMP examples releases the second core)
$\rightarrow$ the runtime's startup, which calls \texttt{adainit} (elaboration)
$\rightarrow$ your \texttt{Main}. There is no vendor startup shim and no
third-party RTOS scheduler --- the Ada runtime is the only scheduler, and it is
linked from source.

\begin{figure}[ht]
\centering
\resizebox{\linewidth}{!}{%
\begin{tikzpicture}[node distance=7mm]
\node[blk] (rom)   {ROM};
\node[blk, right=of rom]   (boot)  {2nd-stage\\bootloader\\\scriptsize cache, PSRAM};
\node[blk, right=of boot]  (start) {\texttt{start.S}\\\scriptsize PLL, vectors,\\\scriptsize core 1};
\node[blk, right=of start] (init)  {\texttt{adainit}\\\scriptsize elaboration};
\node[blk, right=of init]  (main)  {\texttt{Main}};
\draw[flow] (rom)--(boot);   \draw[flow] (boot)--(start);
\draw[flow] (start)--(init); \draw[flow] (init)--(main);
\end{tikzpicture}}
\caption{Reset to \texttt{Main}.}
\label{fig:boot}
\end{figure}

\section{A minimal application}

\begin{lstlisting}
with Ada.Real_Time; use Ada.Real_Time;
with ESP32S3.GPIO;

procedure Main is
   Led : constant ESP32S3.GPIO.Pin_Id := 2;
begin
   ESP32S3.GPIO.Configure (Led, Mode => ESP32S3.GPIO.Output);
   loop
      ESP32S3.GPIO.Toggle (Led);
      delay until Clock + Milliseconds (250);
   end loop;
end Main;
\end{lstlisting}

\texttt{delay until} is served by the board-support tick; \texttt{ESP32S3.GPIO} is
the HAL; everything links against the runtime crate.

\chapter{The Runtime Profiles}
\label{ch:profiles}

The runtime is offered in three \emph{profiles}, selected by the environment
variable \texttt{ESP32S3\_RTS\_PROFILE} when the crate is built. They trade
language features against footprint and certifiability. The same application
source can often build under more than one; the HAL drivers that use finalization
are restricted to the richer two.

\begin{longtable}{p{2.6cm} p{11.3cm}}
\textbf{Profile} & \textbf{What it gives you} \\ \hline
\endhead
\texttt{light-tasking} &
  The Ravenscar/Jorvik subset (\texttt{pragma Profile (Jorvik)}). Tasks,
  protected objects, \texttt{delay until}, priorities --- but with restrictions
  the compiler enforces: \texttt{No\_Finalization}, \texttt{No\_Exception\_%
  Propagation}, \texttt{No\_Secondary\_Stack}, and a heap-less bump allocator.
  Smallest and most analysable; a raised exception is fatal (resets the board).
  HAL drivers that need none of those (GPIO, RNG, Temperature, SDMMC, Touch,
  RTC, crypto) build here. \\
\texttt{embedded} &
  Still Jorvik, but exception-capable: full exception propagation with messages,
  controlled-type \emph{finalization} (RAII), a secondary stack, and a real
  heap (the O(1) TLSF allocator, \S\ref{ch:heap}). This is the default target for most of the HAL and for the
  ext4 filesystem. The RAII driver handles (SPI/I2C/UART sessions, GDMA/LEDC/RMT/
  \dots\ channels) live here. \\
\texttt{full} &
  Complete Ada tasking with no Jorvik restrictions: dynamic task creation, task
  hierarchies, selective \texttt{accept}/rendezvous, \texttt{abort} and
  asynchronous transfer of control. Built on a fully ported GNARL kernel.
  Experimental (work in progress), but real rendezvous and \texttt{abort} run on
  the silicon. \\
\end{longtable}

\section{Choosing a profile}

\begin{itemize}
\item Certifiable, deterministic, minimal-footprint firmware, or code that must
  never allocate: \texttt{light-tasking}.
\item Most applications --- you want exceptions, RAII handles, the secondary
  stack, a heap, and the full HAL including the filesystem: \texttt{embedded}.
\item Code that genuinely needs dynamic tasks, rendezvous or \texttt{abort}:
  \texttt{full}.
\end{itemize}

The selection happens at build time. An example's \texttt{build.sh} sets it:

\begin{lstlisting}[style=sh]
export ESP32S3_RTS_PROFILE=embedded   # or light-tasking / full
\end{lstlisting}

Because the HAL keys its object directory by profile, the same shared driver tree
is safely reused across profiles without stale objects.

\chapter{Building and Running on the ESP32-S3}
\label{ch:build}

\section{Prerequisites}

\begin{itemize}
\item The Alire-managed GNAT cross-toolchain for the target
  (\texttt{gnat\_xtensa\_esp32\_elf}) and \texttt{gprbuild}; these are fetched by
  Alire and discovered automatically by the build scripts.
\item A USB cable to the board's built-in USB-serial-JTAG port (it appears as
  \texttt{/dev/ttyACM0}).
\item Nothing else: no external SDK or build system. An Ada \texttt{esp\_flash}
  host tool (built with the Alire \texttt{gnat\_native} toolchain) packages and
  writes the image.
\end{itemize}

\section{The \texttt{./x} helper}

From the repository root, the \texttt{./x} script builds, flashes, monitors and
debugs any example by name:

\begin{lstlisting}[style=sh]
./x build  esp32s3_gpio0_blink            # cross-compile + link + package app.bin
./x flash  esp32s3_gpio0_blink -p /dev/ttyACM0
./x run    esp32s3_gpio0_blink -p /dev/ttyACM0   # build + flash + serial monitor
\end{lstlisting}

\section{What a build does}

\texttt{build.sh} sets the profile and heap size, then \texttt{bare\_build.sh}:

\begin{enumerate}
\item regenerates/selects the runtime crate for the chosen profile
  (\texttt{gen\_runtime.sh}; how that works, and how to port it, is
  Chapter~\ref{ch:runtime-build});
\item compiles the Ada application against the pinned runtime with
  \texttt{gprbuild}, producing a relocatable \texttt{ada\_app.o}
  (\texttt{Main} renamed to \texttt{\_ada\_main});
\item compiles the shared bare boot (\texttt{start.S}, the interrupt vectors,
  the heap and a freestanding libc) and the example's C console glue;
\item links everything against the linker scripts into \texttt{app.elf}; and
\item packages \texttt{app.bin} (image header + segments).
\end{enumerate}

\texttt{flash.sh} writes the 2nd-stage bootloader, the partition table and
\texttt{app.bin} over USB-serial.

The last two steps --- turning the linked \texttt{app.elf} into a flashable
\texttt{app.bin}, and writing it to the chip --- are handled by two small
\emph{host} tools written in native Ada (built with the Alire
\texttt{gnat\_native} toolchain), so the whole build/flash path needs only the
Alire toolchain and the committed vendored blobs.

\section{Packaging and flashing: two Ada host tools}

\subsection{\texttt{esp\_elf2image} --- ELF to flash image}

\texttt{examples/common/bare/elf2image} converts \texttt{app.elf} into the
ESP32-S3 flash image \texttt{app.bin} --- the exact on-flash format the chip's
ROM and our second-stage bootloader expect (flash mode \texttt{dio}, freq
\texttt{80m}, size \texttt{2MB}). The generated image boots unmodified on
hardware across every example. The packaging is more than a flat copy:

\begin{itemize}
  \item \textbf{Segments} are the ELF's allocated \texttt{PROGBITS} sections,
        merged where contiguous and each padded to a multiple of four.
  \item A \textbf{24-byte header} --- the 8-byte common part (\texttt{magic 16\#E9\#},
        segment count, the \texttt{dio}/\texttt{2MB}/\texttt{80m} flash byte, the
        entry point) plus a 16-byte extended part (\texttt{chip\_id=9},
        \texttt{wp\_pin=16\#EE\#}, \texttt{hash\_appended=1}).
  \item \textbf{Flash (IROM/DROM) segments aligned to 64\,KB}, with the RAM
        segments written \emph{interleaved as the alignment padding} (then a zero
        \texttt{PADDING} segment for the remainder) --- the exact scheme the
        bootloader's MMU mapping relies on.
  \item An \textbf{XOR checksum} (seed \texttt{16\#EF\#}) as the last byte of a
        16-aligned block, and finally a \textbf{SHA-256} of the whole image
        appended (computed by an in-tree \texttt{sha256.adb}).
\end{itemize}

The flash parameters (\texttt{dio}/\texttt{80m}/\texttt{2MB}, chip = ESP32-S3) are
fixed to match the bare boot; they are constants at the top of the source.

\subsection{\texttt{esp\_flash} --- writing it to the chip}

\texttt{examples/common/bare/espflash} writes the image over the chip's
\emph{ROM} serial bootloader.
It is 100\,\% Ada: the OS serial interface (open, termios raw mode, the
\texttt{TIOCM*} DTR/RTS modem lines, poll, read, write) is bound directly to libc
through \texttt{Interfaces.C} --- there is no C source. (\texttt{GNAT.Serial\_-
Communications} would cover the byte I/O but not the DTR/RTS control the reset
needs, so the tool binds \texttt{ioctl(TIOCMBIS/TIOCMBIC)} itself.) The sequence
is the standard ROM protocol, with no RAM stub and no compression:

\begin{enumerate}
  \item \textbf{Reset into download mode} --- the USB-JTAG DTR/RTS pulse.
  \item \textbf{SYNC} --- the \texttt{16\#07\# 16\#07\# 16\#12\# 16\#20\#} + 32$\times$\texttt{16\#55\#}
        handshake until the ROM answers.
  \item \textbf{SPI attach} and \textbf{set-params} (the flash geometry).
  \item Per file, \textbf{flash\_begin} (erase) then \textbf{flash\_data} in
        \texttt{16\#400\#}-byte blocks (the last padded with \texttt{16\#FF\#}), each
        carrying the ROM XOR checksum (seed \texttt{16\#EF\#}).
  \item \textbf{flash\_end}, then a \textbf{hard reset to run} (suppressible with
        \texttt{--no-reset} for capture workflows like the ACATS sweep, which
        does its own reset).
\end{enumerate}

Every command is a \texttt{<direction, op, length, checksum>} packet in SLIP
framing (\texttt{16\#C0\#} delimiters, \texttt{16\#DB\#} escaping) over the serial port.
A typical invocation writes the three images at their flash offsets:

\begin{lstlisting}[style=sh]
esp_flash -p /dev/ttyACM0 --flash-size 2MB \
  0x0     bootloader.bin \
  0x8000  partition-table.bin \
  0x10000 app.bin
\end{lstlisting}

Note that \texttt{esp\_flash} writes \emph{flash} only; the external PSRAM is not
flashed --- the 2nd-stage bootloader maps it at runtime. On hardware a $\sim$225\,KB
image flashes in about three seconds (ROM speed, no stub).

\section{Starting your own project}

The \texttt{esp32-ada} launcher scaffolds a fresh application in any empty
directory (no runtime sources are copied --- it pins the crate):

\begin{lstlisting}[style=sh]
. export.sh                 # put esp32-ada on PATH
mkdir my_app && cd my_app
esp32-ada new               # scaffold build.sh, the .gpr, src/main.adb, board.ads
# edit src/main.adb ...
esp32-ada run -p /dev/ttyACM0
\end{lstlisting}

Each project owns a \texttt{config/board.ads} that states the flash and PSRAM
sizes (used to size the image header and build/select the bootloader):

\begin{lstlisting}
package Board is
   Flash_Size : constant := 2 * 1024 * 1024;   --  2 MB
   PSRAM_Size : constant := 2 * 1024 * 1024;   --  mapped at 16#3D00_0000#
end Board;
\end{lstlisting}

\section{Console output}

The runtime wires \texttt{Ada.Text\_IO} to the built-in USB-serial-JTAG console,
so \texttt{Put\_Line} reaches \texttt{/dev/ttyACM0}. The examples additionally use
a tiny C glue layer calling the ROM \texttt{esp\_rom\_printf} for formatted status
lines. Either way, \texttt{./x run} opens a monitor so you see the output live.
